\section{Conclusion}
\label{sec:conclusion}

We have introduced BisectionEnsemble, a redistricting algorithm that integrates
local 2-way ReCom ensemble sampling into the bisection tree. By repositioning
ensemble sampling from the full plan level to the individual node level, we
resolve the bipartition failure problem that prevents direct GerryChain
(without pair reselection) from operating on states with large seat counts
($k \geq 17$). Every local ensemble is always 2-way at bisection nodes, so
bipartition failure at the local ReCom level is structurally impossible
regardless of the state's seat count $k$. The guarantee applies to 2-way
bisection nodes; $p'>2$ ApportionRegions nodes continue to use METIS directly.

The algorithm is implemented as \texttt{SeedCompositor::BisectionEnsemble\{p,
ensemble\_steps\}} in \texttt{redist-cli}, using
\texttt{split\_subgraph\_bisection\_ensemble()} which seeds a METIS bisection,
runs \texttt{ensemble\_steps} ReCom steps via
\texttt{redist\_ensemble::recom::RecomChain}, and returns the plan at rank
$\lfloor p \times \text{accepted\_count} \rfloor$.

Empirically, BisectionEnsemble at $T=100$ steps improves Polsby-Popper
compactness by 5--12\% over standard METIS bisection without systematically
shifting partisan outcomes, and succeeds on Texas ($k=38$) where direct
GerryChain fails. Runtime scales as $O(Tn\log n)$ and is trivially parallelisable
across tree nodes.

\paragraph{Future work.}
Several directions merit investigation. First, the ReCom chain's mixing time
on local subgraphs is unexplored: how many steps $T$ are needed for convergence
to the stationary distribution? Second, BisectionEnsemble currently applies only
to bisection nodes ($p' = 2$); extending it to $p'$-way nodes with $p' > 2$
(via nested bisections or $p'$-way recombination) would enable full ensemble
coverage of all ApportionRegions tree nodes. Third, the sensitivity of compactness
gains to the ReCom proposal distribution (e.g., biased spanning trees favouring
compact cuts) is a natural direction for further improvement. Finally, a formal
ergodicity proof for the local chain over the feasible bisection space of each
node would strengthen the theoretical foundation.
